\chapter{Heisenberg Equation}

\section*{Introduction}

The Heisenberg equation is $d\hat{A}/dt = (i/\hbar)[\hat{H}, \hat{A}] + \partial\hat{A}/\partial t$. The first term describes the unitary time evolution generated by the Hamiltonian; the second term accounts for explicit time dependence of the operator. For the position and momentum operators, the equation gives $d\hat{x}/dt = \hat{p}/m$ (recovering the classical relation) and $d\hat{p}/dt = -\partial\hat{V}/\partial\hat{x}$ (recovering Newton's second law). This is Ehrenfest's theorem: the expectation values of quantum operators obey classical equations of motion. The commutator $[\hat{H}, \hat{A}]$ encodes how the observable $\hat{A}$ is ``connected'' to the energy of the system. If $\hat{A}$ commutes with $\hat{H}$, it is a constant of motion: $\dot{\hat{A}} = 0$.
\section*{Fundamental Equation Used}

\[
\frac{d\hat{A}}{dt}=\frac{i}{\hbar}[\hat{H},\hat{A}]+\frac{\partial\hat{A}}{\partial t}
\]
\section*{Assumptions}

\textbf{Assumptions:} The Hamiltonian $\hat{H}$ generates time evolution via the unitary operator $\hat{U}(t) = e^{-i\hat{H}t/\hbar}$. Operators are Hermitian (corresponding to observables). The Hilbert space is complete and the operators are densely defined.


\textbf{Limitations:} Does not describe measurement itself (the ``measurement problem'' remains). For open quantum systems coupled to an environment, the equation must be replaced by a master equation (Lindblad form). In relativistic quantum field theory, operators at spacelike separations commute, and the equation must be formulated covariantly.


\section*{Mathematical Derivation}

\textbf{Step 1: The Schr\"odinger picture and time evolution.}

In the Schr\"odinger picture, the state vector $|\psi_S(t)\rangle$ evolves according to the Schr\"odinger equation:
\begin{align}
i\hbar\frac{d}{dt}|\psi_S(t)\rangle = \hat{H}|\psi_S(t)\rangle,
\end{align}
with formal solution $|\psi_S(t)\rangle = \hat{U}(t)|\psi_S(0)\rangle$, where $\hat{U}(t) = e^{-i\hat{H}t/\hbar}$ is the time-evolution operator (for time-independent $\hat{H}$). Operators $\hat{A}_S$ are time-independent in this picture.

\textbf{Step 2: Transition to the Heisenberg picture.}

Define the Heisenberg-picture state as $|\psi_H\rangle = |\psi_S(0)\rangle$ (fixed, time-independent) and the Heisenberg operator as
\begin{align}
\hat{A}_H(t) = \hat{U}^\dagger(t)\,\hat{A}_S\,\hat{U}(t) = e^{i\hat{H}t/\hbar}\,\hat{A}_S\,e^{-i\hat{H}t/\hbar}.
\end{align}
Expectation values are the same in both pictures: $\langle\psi_S(t)|\hat{A}_S|\psi_S(t)\rangle = \langle\psi_H|\hat{A}_H(t)|\psi_H\rangle$.

\textbf{Step 3: Differentiate to find the equation of motion.}

Taking the time derivative of $\hat{A}_H(t)$:
\begin{align}
\frac{d\hat{A}_H}{dt} &= \frac{i}{\hbar}\hat{H}\,e^{i\hat{H}t/\hbar}\hat{A}_S\,e^{-i\hat{H}t/\hbar} + e^{i\hat{H}t/\hbar}\hat{A}_S\left(-\frac{i}{\hbar}\hat{H}\right)e^{-i\hat{H}t/\hbar} \\
&= \frac{i}{\hbar}\bigl(\hat{H}\hat{A}_H(t) - \hat{A}_H(t)\hat{H}\bigr) \\
&= \frac{i}{\hbar}[\hat{H}, \hat{A}_H(t)].
\end{align}

\textbf{Step 4: Include explicit time dependence.}

If $\hat{A}_S$ depends explicitly on time (e.g., an operator with time-dependent parameters), the full Heisenberg equation is:
\begin{align}
\boxed{\frac{d\hat{A}_H}{dt} = \frac{i}{\hbar}[\hat{H}, \hat{A}_H] + \frac{\partial\hat{A}_H}{\partial t},}
\end{align}
where the last term is the Heisenberg-picture derivative of the explicit time dependence. For most standard operators (position, momentum, spin), $\partial\hat{A}_H/\partial t = 0$.

\textbf{Step 5: Ehrenfest's theorem.}

For $\hat{A} = \hat{x}$ and $\hat{H} = \hat{p}^2/(2m) + V(\hat{x})$, the commutators are $[\hat{H}, \hat{x}] = -i\hbar\hat{p}/m$ and $[\hat{H}, \hat{p}] = -i\hbar\,dV/d\hat{x}$ (using $[\hat{p}, V] = -i\hbar\,dV/d\hat{x}$). The Heisenberg equations become:
\begin{align}
\frac{d\hat{x}}{dt} = \frac{\hat{p}}{m}, \qquad \frac{d\hat{p}}{dt} = -\frac{dV}{d\hat{x}},
\end{align}
which are the operator forms of Newton's second law. Taking expectation values yields Ehrenfest's theorem: $m\,d\langle x\rangle/dt = \langle p\rangle$ and $d\langle p\rangle/dt = -\langle dV/dx\rangle$.

\section*{Characteristics of the Equation}

The Heisenberg equation preserves the algebra of observables: if two operators satisfy a commutation relation at $t = 0$, the evolved operators satisfy the same relation at all times. This is because unitary evolution is an automorphism of the algebra. The equation is linear: the evolution of $\hat{A} + \hat{B}$ is the sum of the evolutions. It is also covariant under unitary transformations: if $\hat{A} \to \hat{U}\hat{A}\hat{U}^\dagger$, the equation retains its form. For time-independent operators, $d\hat{A}/dt = (i/\hbar)[\hat{H}, \hat{A}]$, which is a purely algebraic equation.
\section*{Solution Methods}

\textbf{Direct computation:} Compute the commutator $[\hat{H}, \hat{A}]$ using the known algebra of operators. For the harmonic oscillator, $[\hat{H}, \hat{a}] = -\hbar\omega\hat{a}$, giving $\hat{a}(t) = \hat{a}(0)e^{-i\omega t}$. \textbf{Heisenberg operators:} Express Heisenberg-picture operators as $\hat{A}_H(t) = \hat{U}^\dagger(t)\hat{A}_S\hat{U}(t)$ and expand. \textbf{Interaction picture:} Split $\hat{H} = \hat{H}_0 + \hat{V}$ and work in the interaction picture where $\hat{H}_0$ evolution is removed. \textbf{Numerical methods:} For complex systems, solve the Heisenberg equation numerically using matrix exponentiation or Trotter decomposition.
\section*{Verifications}

For a free particle ($\hat{H} = \hat{p}^2/2m$), check that $d\hat{x}/dt = \hat{p}/m$ and $d\hat{p}/dt = 0$. For the harmonic oscillator ($\hat{H} = \hbar\omega(\hat{a}^\dagger\hat{a} + 1/2)$), verify that $\hat{a}(t) = \hat{a}(0)e^{-i\omega t}$. Check Ehrenfest's theorem: $\frac{d}{dt}\langle\hat{x}\rangle = \langle\hat{p}\rangle/m$ and $\frac{d}{dt}\langle\hat{p}\rangle = -\langle\partial V/\partial x\rangle$. Verify that the Heisenberg uncertainty relation is preserved under time evolution. Check that operators commuting with $\hat{H}$ are constants of motion.
\section*{Applications}

The Heisenberg equation is the foundation of quantum field theory, where fields are operators that evolve according to the Heisenberg equation with a field-theoretic Hamiltonian. In condensed matter physics, it describes the dynamics of spin operators in the Heisenberg model of magnetism. In quantum optics, it governs the evolution of creation and annihilation operators in the Jaynes--Cummings model. In quantum computing, it describes the evolution of qubit operators under Hamiltonian gates. In quantum information, it provides the equations of motion for entanglement entropy and other information-theoretic quantities.
\section*{What Richard Feynman Would Have Asked}

\subsection*{1. Plain-English Translation}
\textit{``What is the Heisenberg equation telling us that the Schr\"odinger equation does not?''}

The Heisenberg equation tells us how observables---the quantities we actually measure---evolve in time. While the Schr\"odinger equation describes how the wave function (an abstract mathematical object) changes, the Heisenberg equation describes how the physical quantities themselves change. It is the quantum version of Newton's second law: instead of $\dot{p} = F$, we have $d\hat{p}/dt = (i/\hbar)[\hat{H}, \hat{p}]$.

The commutator $[\hat{H}, \hat{p}]$ measures how ``different'' the momentum is from the energy. If they commute ($[\hat{H}, \hat{p}] = 0$), the momentum is a constant of motion---it does not change. If they do not commute, the momentum evolves. The rate of change is proportional to the size of the commutator and inversely proportional to $\hbar$: smaller $\hbar$ means faster, more classical-looking evolution.

\subsection*{2. Localized Mechanics}
\textit{``At a single instant, what determines how fast an observable changes?''}

At any instant, the rate of change of an observable $\hat{A}$ is determined by its commutator with the Hamiltonian: $[\hat{H}, \hat{A}]$. This commutator measures the extent to which $\hat{A}$ ``interferes'' with the energy. If $\hat{A}$ is a conserved quantity (like total energy, $[\hat{H}, \hat{H}] = 0$), it does not change. If $\hat{A}$ is the position, the commutator $[\hat{H}, \hat{x}] = -i\hbar\hat{p}/m$ gives $d\hat{x}/dt = \hat{p}/m$---the familiar velocity-momentum relation.

The commutator also encodes the uncertainty principle. If $[\hat{A}, \hat{B}] = i\hbar\hat{C}$, then $\hat{A}$ and $\hat{B}$ cannot both be known precisely at the same time. The Heisenberg equation shows how this uncertainty evolves: even if you prepare a state with precise values of certain observables, the commutator with $\hat{H}$ can generate uncertainty over time.

\subsection*{3. Testing Extreme Limits}
\textit{``What happens when $\hbar \to 0$, when the system is strongly coupled, or at extreme energies?''}

Classical Limit ($\hbar \to 0$): The commutator $[\hat{H}, \hat{A}]$ becomes $i\hbar\{H, A\}$ (the classical Poisson bracket), and the Heisenberg equation reduces to $\dot{A} = \{A, H\}$---the classical equation of motion. The quantum corrections are of order $\hbar$ and vanish in the classical limit. This is why classical mechanics works so well for macroscopic objects.

Strong Coupling: When the Hamiltonian has strong off-diagonal elements (strong coupling between states), the Heisenberg equation produces rapid oscillations. In quantum optics, this leads to Rabi oscillations between two states. In condensed matter physics, it leads to spin waves and magnons.

Extreme Energies: At energies where particle creation becomes important (quantum field theory), the Heisenberg equation for fields produces particle-antiparticle pair creation. The Hamiltonian includes terms that create and destroy particles, and the Heisenberg equation governs their production rates.

\subsection*{4. Dimensionality and Geometry}
\textit{``Does the Heisenberg equation depend on the dimensionality of the Hilbert space?''}

The Heisenberg equation is algebraic---it depends only on the commutation relations, not on the dimensionality of the Hilbert space. In a 2D Hilbert space (a single qubit), the operators are $2\times 2$ matrices and the Heisenberg equation is a matrix ODE. In an infinite-dimensional Hilbert space (a harmonic oscillator), the operators are infinite matrices and the equation is still the same. The physics (unitary evolution, preservation of commutation relations) is identical regardless of dimension.

In quantum field theory, the Hilbert space is infinite-dimensional and the operators are fields. The Heisenberg equation becomes a field equation---the quantum equations of motion for the fields. The geometry of spacetime enters through the Hamiltonian (which depends on the metric in curved spacetime) and through the commutation relations (which are local in relativistic theories).

\subsection*{5. Cross-Disciplinary Connection}
\textit{``Where else does the Heisenberg equation appear outside of quantum mechanics?''}

The mathematical structure---operators evolving under a commutator with a ``Hamiltonian''---appears in many fields. In classical statistical mechanics, the Liouville equation for the phase space distribution can be written in a form analogous to the Heisenberg equation. In signal processing, the time evolution of operators in the Wigner function formalism has the same structure. In machine learning, the ``Hamiltonian'' framework has been applied to neural network dynamics, where the ``operators'' are weight matrices and the ``commutator'' governs their evolution.

In control theory, the evolution of the state-transition matrix follows an equation that is mathematically identical to the Heisenberg equation. The deep reason is that any linear, reversible, norm-preserving dynamics can be described by a generator (the ``Hamiltonian'') acting through a commutator. The Heisenberg equation is the most general expression of this mathematical structure.

%=============================================================
\section*{FAQ}

\textbf{Q: What is the difference between the Heisenberg and Schr\"odinger pictures?}

A: In the Schr\"odinger picture, the state vector evolves in time while operators are fixed. In the Heisenberg picture, operators evolve while the state vector is fixed. Both give identical predictions for expectation values: $\langle\psi_S|\hat{A}_S|\psi_S(t)\rangle = \langle\psi_H(t)|\hat{A}_H(t)|\psi_H\rangle$. The Heisenberg picture is more natural for relativistic theories and for connecting to classical mechanics.

\textbf{Q: What does Ehrenfest's theorem really mean?}

A: Ehrenfest's theorem says that the expectation values of position and momentum obey Newton's equations: $m\,d\langle x\rangle/dt = \langle p\rangle$ and $d\langle p\rangle/dt = -\langle\partial V/\partial x\rangle$. The classical equation has $\partial V/\partial x$ evaluated at $x$; the quantum equation has $\langle\partial V/\partial x\rangle$, which is the average of the force over the wave function. For a linear force (harmonic oscillator), the two are equal. For a nonlinear force, they differ---this is the origin of quantum deviations from classical behavior.

\textbf{Q: Can the Heisenberg equation handle dissipation?}

A: Not directly. The Heisenberg equation describes closed, isolated systems. For open quantum systems (coupled to an environment), the evolution is non-unitary and is described by a Lindblad master equation or a quantum Langevin equation. These are generalizations of the Heisenberg equation that include dissipation and noise.
\section*{Important Bibliography}

\begin{enumerate}
\item Sakurai, J.J. and Napolitano, J., \textit{Modern Quantum Mechanics}, 2nd ed. (Cambridge University Press, 2017), Chapter 2.
\item Cohen-Tannoudji, C., Diu, B., and Laloe, F., \textit{Quantum Mechanics}, 2nd ed. (Wiley-VCH, 2020), Chapter VI.
\item Heisenberg, W., ``Quantentheoretische Umdeutung kinematischer und mechanischer Beziehungen,'' \textit{Zeitschrift f\"ur Physik} \textbf{33}, 879--893 (1925).
\item Messiah, A., \textit{Quantum Mechanics}, 2nd ed. (Dover, 1999), Chapter VII.
\end{enumerate}


